% References and numerical settings checked against the recorded experiments.
% The existing performance runs use Daejeon geometry. The current Yeouido map
% is a separately generated illustration, not the geometry of those runs.
% Restore the Yeouido sentence and map cross-reference only after corresponding
% performance runs have actually used that geometry.
\subsection{Simulation Settings}
\label{subsec:simulation_settings}

In the simulations, we consider $U=10$ UAVs operating over an
$1800\times1800$~m$^2$ area in Daejeon, South Korea, at an altitude
of 50~m to represent an environment with diverse IoT devices
and large volumes of collected data.
The UAVs follow prescribed trajectories and hover during each
communication interval.
For the UAV-to-UAV links, we adopt the free-space path-loss model
used in UAV-assisted AirComp \cite{Zhai_TWC_2025}, with a reference
channel power gain of $-60$~dB at 1~m, receiver noise power
of $-90$~dBm, and maximum transmit power of 0.32~W.
Channels remain constant within each transmission block,
with synchronized transmission and perfect channel state
information assumed.
Unless otherwise specified, the uplink bandwidth is 1~MHz,
the block coherence time is 10~ms, and the network RF
transmission energy budget per round is 0.1~J.

The computation delay and dynamic CPU energy are modeled
using the executed CPU cycles and processor frequency
\cite{Tang_TWC_2025,Liang_TWC_2025_Energy}.
We set the processor frequency to 1.5~GHz, the effective
switched capacitance to $10^{-27}$, and the active board
power to 2~W for each UAV.
The computation workload is estimated from the executed
multiply--accumulate operations, assuming eight operations
per CPU cycle.
The hovering power is 168.49~W, corresponding to the sum
of the blade-profile and induced powers of 79.86 and
88.63~W, respectively \cite{Wang_TCOMM_2025_Rechargeable}.
The power amplifier efficiency is set to 0.35.

The UAVs perform video--text retrieval using CapERA as a UAV
video-caption benchmark \cite{Bashmal_RS_2023_CapERA} and a
WebUAV3M-source subset of AirScape as a UAV video-intention
benchmark \cite{Zhao_ACMMM_2025_AirScape,Zhang_TPAMI_2023_WebUAV3M}.
Training samples are distributed among the UAVs using a label-based
Dirichlet partition with $\alpha_{\mathrm D}=0.1$
\cite{Mora_TBD_2026,Liang_TWC_2025_Energy}.
We use pretrained RemoteCLIP with a ViT-B/32
visual backbone \cite{Liu_TGRS_2024_RemoteCLIP}, freeze both
encoders and the shared logit-scale parameter, and train rank-10
residual adapters for the visual and textual features,
yielding 20,480 trainable parameters.
The initial learning rate is $3.5\times10^{-3}$.
All UAVs participate in each feasible FL round and perform
one SGD step without momentum, using a mini-batch size of 16.
The learning rate remains constant through round 40 and follows a cosine
decay over rounds 41--200, and the same rate is used for unlearning.

Each run comprises 20 initial learning rounds followed by 180 rounds
with unlearning requests. Each UAV initially has access to its entire
local dataset, with no subsequent data arrivals. Five UAVs are uniformly
selected from eligible UAVs, each requesting the deletion of half of
the samples used during the initial learning rounds. The selected
samples are divided into nine unlearning requests per UAV, which are
issued at predetermined rounds for each dataset.
The deletion fraction follows~\cite{Mora_TBD_2026}, while the request
schedules are selected for this study. Requested samples are unlearned
over five update steps per request and excluded from subsequent learning
upon request arrival.
The following benchmark schemes are used to evaluate the performance
of the proposed framework.
